\begin{proof}

\begin{enumerate}[label={$\langle1\rangle$\arabic*.},leftmargin=*]
\item Restriction of the crossed homomorphism to the subgroup generated by
$\rho_j^{m_j}$ has value
\[
 z(\rho_j^{m_j})
 =\sum_{r=0}^{m_j-1}T_j^rz(\rho_j)=N_{T_j}d_j.
\tag{W.29}\hypertarget{tag-03GW}{}\label{wbII:eq:W:5.17l}
\]
This is the cohomological pullback of the annular defect to the
$s_j$-cover; the norm, rather than unqualified multiplication by $m_j$, is
essential before passage to global coinvariants.  A principal change
$d_j\mapsto d_j+(T_j-I)a$ does not alter this value, since
$N_{T_j}(T_j-I)a=(T_j^{m_j}-I)a=0$.

\item Let $Q_j:J'_j\to N_j$ be the finite quotient map.  The global flat marking
$\Xi$ over the full disc $\Delta_j$ makes the constant vertical form
$\Omega_E$ well defined on $J'_j$.  Radial retraction onto the central
torus makes
\[
 H^2(J'_j;\mathbb R)\longrightarrow H^2(A_0^{(j)};\mathbb R)
\]
an isomorphism.  On the central fibre the quotient is exactly $\pi_j$, and
naturality gives
\[
 (Q_j|_{A_0^{(j)}})^*c_1(L|_{S_j})
 =\pi_j^*c_1(L|_{S_j})=E.
\tag{W.30}\hypertarget{tag-03GX}{}\label{wbII:eq:W:5.17l'}
\]
The second equality is the actual radial/stalk identification in
\CompendiumRef{I}{wb:lem:finite-filling-stalk-map}; on a nearby fibre the gluing map $h_j$
is a translation and therefore acts trivially on $H^2$ of that fibre.
Consequently the two classes
$c_1(Q_j^*(L|_{N_j}))$ and $[\Omega_E]$ have the same central restriction,
and the retraction isomorphism proves
\[
 c_1(Q_j^*(L|_{N_j}))-[\Omega_E]=0
 \quad\text{in }H^2(J'_j;\mathbb R).
\tag{W.31}\hypertarget{tag-03GY}{}\label{wbII:eq:W:5.17l''}
\]

\item On $\Delta_j^*$, the source identity
$h_j\widetilde g_j^{\log}h_j^{-1}=\widetilde g_j^0$ shows that $h_j$
descends from the logarithmic quotient $N_j^*$ to the linear/global
quotient $J^*$.  Thus
\[
 Q_j^*(L|_{N_j})|_{\Delta_j^*}
 \simeq h_j^*(L|_{J^*}).
\tag{W.32}\hypertarget{tag-03GZ}{}\label{wbII:eq:W:5.17l'''}
\]
The logarithmic overlap is the holomorphic translation
\[
 h_j[s,\zeta]=[s,\zeta+\alpha(s)\Pi(s)v_j],
 \qquad \alpha(s)=\frac{\log s_j}{2\pi i}.
\]
To compute its effect on a real cohomology class, retain the exact
real marking
\[
 \Xi(s,x+\Lambda)=[s,\Pi(s)x],\qquad
 J_s=\Pi(s)^{-1}\circ i\circ\Pi(s)
 \quad\text{on }\Lambda_{\mathbb R}.
\]
The period matrix \(\Pi(s)\) is a real isomorphism by the displayed period inequality;
thus \(J_s\) is a well-defined smooth family of real complex
structures.  On a branch of the logarithm the actual real
translation vector is
\[
 a_j(s)=\operatorname{Re}\alpha(s)\,v_j+
        \operatorname{Im}\alpha(s)\,J_s v_j.
 \tag{W.33}\hypertarget{tag-03H0}{}\label{wbII:eq:W:5.17m}
\]
Indeed, applying \(\Pi(s)\) gives precisely
\(\alpha(s)\Pi(s)v_j\).  A different logarithm branch changes
\(a_j\) by an integral multiple of \(v_j\), so its torus translation
is branch-independent.  Complex multiplication by \(\alpha\) is
not multiplication by a real scalar in this marking.

\item Write \(E=\frac12\sum_{a,b}E_{ab}\,dx^a\wedge dx^b\) in the
ordered real coordinates dual to the marked lattice; the matrix
\((E_{ab})\) is alternating and is constant over the covering
disc.  Direct substitution in the real map
\(h_j(s,x)=(s,x+a_j(s))\) gives
\[
 \begin{split}
 h_j^*\Omega_E-\Omega_E
  ={}&\sum_{a,b}E_{ab}\,d a_j^a\wedge dx^b\\
    &+\frac12\sum_{a,b}E_{ab}\,d a_j^a\wedge d a_j^b.
 \end{split}
 \tag{W.34}\hypertarget{tag-03H1}{}\label{wbII:eq:W:5.17m'}
\]
The last term has both arguments in the base directions.
In particular, it restricts to zero when the base is an oriented
circle.

\item Take one full clockwise circle in the \(s_j\)-cover,
\[
 s_j(\theta)=r e^{-2\pi i\theta},\qquad 0\le\theta\le1.
\]
Choose the logarithm continuously along that circle, so
\(\alpha(1)=\alpha(0)-1\).  The endpoints represent the same
point of the covering disc.  The marking is defined on the full
disc, and therefore \(J_{s_j(1)}=J_{s_j(0)}\).  Consequently
\[
 a_j(1)-a_j(0)=-v_j.
 \tag{W.35}\hypertarget{tag-03H2}{}\label{wbII:eq:W:5.17m''}
\]
For any \(\lambda\in\Lambda\), restrict \eqref{wbII:eq:W:5.17m'} to the sweep of
the marked circle \(x(u)=u\lambda\), using the base-first
orientation \(d\theta\wedge du\).  Its integral is exactly
\[
 \int_0^1 E\!\left(\frac{d a_j}{d\theta},\lambda\right)d\theta
 =E(a_j(1)-a_j(0),\lambda)
 =-E(v_j,\lambda).
\]
The horizontal term contributes zero.  Since marked lattice
circles span \(H_1(F;\mathbb R)\), these integrals determine the
mixed component in
\[
 H^1(S^1;\mathbb R)\otimes H^1(F;\mathbb R).
\]
Thus the mixed component of \(h_j^*\Omega_E-\Omega_E\) is
\(-\iota_{v_j}E\), with \((\iota_{v_j}E)(\lambda)=E(v_j,\lambda)\).

\item The pullback of the defect
\(c_1(L|_{J^*})-[\Omega_E]\) has zero fibre component and mixed
component \(N_{T_j}d_j\) by \eqref{wbII:eq:W:5.17l}.  Translation by \(h_j\) does
not change this mixed component: on the covering circle such
a fibre-zero class is represented by
\(d\theta\wedge\beta\), with \(\beta\) a constant fibre one-form,
and substituting \(dx\mapsto dx+d a_j\) adds only a base two-form,
which vanishes on the circle.  Equations \eqref{wbII:eq:W:5.17l''} and
\eqref{wbII:eq:W:5.17l'''} state that the total pulled-back class extends over
the covering disc and has zero mixed component.  It follows that
\[
 N_{T_j}d_j-\iota_{v_j}E=0.
\]
In the global coinvariants every \(T_j\) acts as the identity,
so \(m_j\bar d_j=[\iota_{v_j}E]\).  This proves \eqref{wbII:eq:W:5.17k}, with the
original clockwise convention and no change to either translation
vector.

\end{enumerate}
\end{proof}
